\begin{description}
\item[(OP02.1)]
% The solution is the marked sky map: the positions of the 10 moons are indicated by
% scoring circles on the solution overlay.
% FIGURE NEEDED: solution overlay of Map-OP02 with the moons' positions circled (inner and
% outer tolerance circles) (2025_Map-OP02-Solutions.pdf).
For each moon:
\begin{itemize}
\item Centre of the $\otimes$ within the inner circle \hfill 2.0
\item Centre of the $\otimes$ between the inner and outer circles \hfill 1.0
\end{itemize}
Use 4 correctly placed moons for marking. Every incorrectly placed moon beyond 4 gets
$-1.0$, upto a minimum of 0. % sic: "upto"

\item[(OP02.2)]
The different angular extents of the shadow on the inner and outer edges of the ring are
visible on the dome. These edges may be marked on the sky map (using convenient positions
of the stars nearby).

From the grid of the RA lines on the map the angular sizes are measured to be:
\begin{center}
\begin{tabular}{ll}
Inner edge: & $\approx 75^\circ$\\
Outer edge: & $\approx 33^\circ$\\
\end{tabular}
\end{center}
These angles need to be used in the schematic diagram to locate the rings.

Place a protractor with its baseline along the dotted line, and with its origin at P (the
position of the observer).

Find the line of sight to the inner edge that subtends an angle of $75^\circ/2 = 37.5^\circ$
at P. Extend this line to intersect the tangent ray. This point marks one end of the shadow
on the inner ring. Similarly, find the other ends of the shadow on the inner edge.

Draw a circular arc (with centre at the centre of the planet) passing through these two
points of intersection to show the inner edge of the ring. Repeat for the outer edge.
% FIGURE NEEDED: completed construction on the polar-view diagram (2025_PLAN_sol.pdf
% pp.5-6): sight lines from P at half the measured angles on either side of the dotted
% line, intersection points I1, I2 (inner edge) and O1, O2 (outer edge) on the tangent
% light rays, and the two circular arcs centred on the planet's centre through those
% points.

\item[(OP02.3)]
Once the edges of the ring are drawn, the width can be calculated by measuring the
difference of their distances from the centre of the planet. Use a ruler to measure these
distances. Also measure the radius of the planet on the diagram with the ruler to find the
calibration.

We find,
\begin{itemize}
\item Distance between the two rings along the dotted line: 3.8 cm
\item Radius of the planet on the diagram: 2 cm equivalent to $6.00\times10^4$ km (given)
\end{itemize}
Therefore,
\[
\text{Width of the ring in km:}
= \frac{3.8\,\mathrm{cm}}{2\,\mathrm{cm}}\times 6.00\times10^4\,\mathrm{km}
= 11.4\times10^4\,\mathrm{km}
\]

Tolerance on values of radii of inner and outer edges:
\begin{center}
\begin{tabular}{|c|c|c|}
\hline
Half credit lower limit & Full credit range & Half credit upper limit\\
\hline
$9.45\times10^4$ km & $10.60\times10^4$ km to $13.40\times10^4$ km & $14.55\times10^4$ km\\
\hline
\end{tabular}
\end{center}
\end{description}
